\documentclass[12pt]{article}

%----------------- Packages ------------------
\usepackage[utf8]{inputenc}
\usepackage[T1]{fontenc}
\usepackage{amsmath, amssymb}
\usepackage{geometry}
\usepackage{xcolor}
\usepackage{titlesec}
\usepackage{fancyhdr}
\usepackage[breakable]{tcolorbox}
\usepackage{graphicx}
\usepackage{tikz}
\usepackage{float}
\usetikzlibrary{arrows.meta, decorations.markings}

%----------------- Page Setup -----------------
\geometry{a4paper, margin=2.5cm}
\pagestyle{fancy}
\fancyhf{}
\rhead{Final Exam — \textbf{2024}}
\lhead{Analysis IV}
\cfoot{\thepage}

%----------------- Title Styling --------------
\titleformat{\section}{\normalfont\Large\bfseries}{Exercise \thesection:}{1em}{}
\titleformat{\subsection}{\normalfont\bfseries}{Answer:}{1em}{}

%----------------- Custom Boxes ----------------
\tcbuselibrary{listingsutf8}
\newtcolorbox{answerbox}{
  colback=gray!10,
  colframe=black,
  fonttitle=\bfseries,
  title=Answer Area,
  breakable,
  before skip=10pt,
  after skip=10pt
}

%----------------- Document Start --------------
\begin{document}

%----------------- Exam Info -------------------
\begin{center}
  \Large\textbf{UNIVERSITY NAME} \\[1em]
  \large\textit{Analysis IV — Final Exam} \\[0.5em]
  \large\textit{Year: 2024} \\[2em]
\end{center}

\vspace{0.5cm}

%----------------- EXERCISE 1 ------------------
\section{}
Determine the nature (convergence/divergence) of the following numerical series of general term:

\[
\begin{aligned}
&a) \; u_n = \frac{n^n}{(2n)!} \\
&b) \; u_n = \frac{n + \cos n}{e^n + \sin n} \\
&c) \; u_n = \sqrt{1 + \frac{(-1)^n}{\sqrt{n}}} - 1
\end{aligned}
\]

\begin{answerbox}

\end{answerbox}

%----------------- EXERCISE 2 ------------------
\section{}
Consider the series of functions \(\sum b_n\) defined by:

\[
f_n(x) = \frac{1}{n + n^2 x}
\]

We denote:

\[
f(x) = \sum_{n=1}^\infty f_n(x)
\]

\begin{enumerate}
    \item Determine the domain of definition of \(f\).
    
    \item Show that \(f\) is continuous on \(\mathbb{R}^*_+\) and compute: 
    $$
    \lim_{x \to \infty} f(x) ; \lim_{x \to 0} f(x)
    $$
    deduction. 
    
    \item Show that \(f\) is of class \(C^1\) on \(\mathbb{R}^*_+\) and give its derivative.
    
    \item Sketch the variation table of \(f\) on \(\mathbb{R}^*_+\) and draw its graph.
    
    \item 
    \begin{enumerate}
      \item Show the existence and compute the integral:
      \[
      \int_1^\infty \frac{1}{y + y^2 x} \, dy
      \]
      
      \item Show that:
      \[
      \log \frac{x+1}{x} \leq f(x) \leq \log \left( \frac{1+2x}{2x} \right) + \frac{1}{1+x}
      \]
      Deduce an equivalent of \( f(x) \) at 0.
    \end{enumerate}   
\end{enumerate}
\begin{answerbox}

\end{answerbox}

%----------------- EXERCISE 3 ------------------
\section{}
Let \( f : x \in (-1, 1) \longrightarrow \frac{\arcsin x}{\sqrt{1-x^2}} \).

\begin{enumerate}
    \item Justify that \( f \) can be expanded as a power series.
    
    \item Show that \( f \) satisfies a first-order differential equation.
    
    \item Deduce the power series expansion of \( f \).
    
    \item Deduce that for any integer \( n \):
    \[
    \sum_{k=0}^{n} \frac{1}{2k+1} \binom{2k}{k} \binom{2(n-k)}{n-k} = \frac{2^{4n}(n!)^2}{(2n+1)!}
    \]
    It is recalled that:
    \[
    \frac{1}{\sqrt{1-x^2}} = \sum_{n=0}^{\infty} \frac{(2n)!}{2^{2n}(n!)^2} x^{2n}
    \]
\end{enumerate}

\begin{answerbox}

\end{answerbox}

%----------------- END ------------------
\end{document}